\chapter{Introduzione e obiettivi}
\section{Problema fisico}
Un manipolatore con attuatori elastici in serie non trasmette rigidamente il moto dei motori ai link. Fra ciascun rotore e il braccio si trova una molla torsionale: il motore deve deformarla per produrre la coppia sul giunto. Questa elasticità rende il comportamento più ricco di quello di un manipolatore rigido e introduce oscillazioni che il controllo deve smorzare.

Studiamo il manipolatore planare a due giunti assegnato nella traccia del corso. All'estremità del secondo link consideriamo anche una massa aggiuntiva, che rappresenta il carico dell'end-effector. Il suo valore può cambiare durante l'impiego del robot. Anche i due drive sono conosciuti solo approssimativamente: frequenza naturale, smorzamento e ritardo possono discostarsi dai valori nominali.

Il punto centrale del lavoro è confrontare due modi di descrivere queste incertezze. Nel modello A manteniamo espliciti i parametri fisici dei drive; nel modello B raccogliamo le loro variazioni in due perturbazioni dinamiche pesate. La massa del carico rimane invece un parametro reale in entrambi. In questo modo il confronto riguarda due descrizioni degli stessi attuatori e della stessa meccanica.

\section{Obiettivi di controllo}
Vogliamo stabilizzare il braccio nella configurazione verticale superiore, seguire piccoli riferimenti articolari e limitare gli effetti di coppie disturbanti applicate ai link. Confrontiamo LQG, LQGI, controllo $H_\infty$ a uno e due gradi di libertà, PID strutturato e sintesi robusta mediante iterazione D--K.

Queste strategie rispondono a domande diverse. LQG bilancia stato e comando e ricostruisce gli stati non misurati; l'integratore corregge gli errori persistenti. La sintesi $H_\infty$ limita una norma pesata, mentre la sintesi con $\mu$ tiene conto della struttura delle incertezze. Per valutarle usiamo gli stessi riferimenti, disturbi e impianti fisici campionati.

\section{Metodo di lavoro e classi di verifica}
Distinguiamo quattro proprietà. La stabilità nominale, NS, riguarda la convergenza degli stati nel modello nominale. La prestazione nominale, NP, richiede anche che una mappa pesata abbia norma minore di uno. La stabilità robusta, RS, estende la stabilità a tutta la famiglia ammessa. La prestazione robusta, RP, richiede stabilità e limite prestazionale per tutti i suoi elementi.

Queste definizioni non sostituiscono le prove temporali. Un controllore può soddisfare una norma pesata e avere un transitorio poco adatto, oppure amplificare il rumore fino a richiedere coppie irragionevoli. Per questo affianchiamo all'analisi robusta tracking, disturbi, sforzo di controllo e simulazioni non lineari. I campioni danno evidenza numerica; non dimostrano da soli una proprietà sull'intero insieme.

\section{Organizzazione della relazione}
Partiamo dalle equazioni fisiche, dall'equilibrio e dalla linearizzazione. Introduciamo poi i drive e le due rappresentazioni delle incertezze. Dopo aver definito obiettivi e pesi, progettiamo i controllori e ne discutiamo la robustezza. Le simulazioni non lineari, la regione di attrazione empirica e i dieci impianti comuni permettono infine di interpretare vantaggi e limiti delle soluzioni. Le appendici raccolgono parametri, istruzioni di riproduzione e pochi estratti MATLAB essenziali.
